\documentclass[10pt,a4paper]{amsart}
\usepackage[margin=19mm]{geometry}
\usepackage{amsmath,amssymb,mathtools}
\usepackage[scr=rsfs,cal=euler]{mathalfa}
\usepackage{xfrac}
\usepackage[unicode,hidelinks,bookmarks=false]{hyperref}
\newcommand{\ie}{i.e.\ }
\newcommand{\cf}{cf.\ }
\newcommand{\resp}{respectively\ }
\newcommand{\Subsection}{\subsection}
\providecommand{\nobreakdash}{\nobreak\hbox{-}}
\setlength{\emergencystretch}{2em}
\multlinegap0pt
\usepackage{amssymb}
\usepackage[shortlabels]{enumitem}
\usepackage{xcolor}
\usepackage{mathtools}
\newtheorem{proposition}{Proposition}
\newcommand{\Grp}{\ensuremath{\mathsf{Grp}}}
\DeclareMathOperator{\bZ}{Z}
\title{Weak action representability\\of \texorpdfstring{\(2\)}{2}-nilpotent groups}
\author{Alessandro Dioguardi Burgio, Manuel Mancini, Tim Van der Linden}
\date{Source: arXiv:2604.22578v1 (CC BY 4.0)}
\begin{document}
\maketitle

\noindent\emph{Source and licence.} Weak action representability\\of \texorpdfstring{\(2\)}{2}-nilpotent groups. Version arXiv:2604.22578v1 (CC BY 4.0), distributed by arXiv under the Creative Commons Attribution 4.0 International licence, http://creativecommons.org/licenses/by/4.0/, as recorded in the arXiv licence metadata for this exact version and shown on the abstract page https://arxiv.org/abs/2604.22578v1. Full licence notice: \texttt{licenses/arxiv-2604.22578-CC-BY-4.0.txt}.\par\smallskip

\noindent\textit{Editorial scope.} Counted unit: the article's proposition PropAct (source0.tex lines 332--349). The article's complete original proof of it is delivered with it (source0.tex lines 351--435, one \texttt{proof} environment of 85 lines). No mathematics is written, repaired or re-derived: the statement and the proof are the article's own lines, in the article's own notation, and no retained line is substituted or re-flowed.  No further source-local context is needed: every cross-reference of every retained line resolves inside the two delivered spans.  Nothing was omitted from any retained span, so the package carries no omission record.  No retained line contains a citation, so the wrapper carries no bibliography and no bibliography entry.  No retained line contains a figure environment, an image-inclusion command or drawing code, so no image is delivered and the package declares no asset dependency.  Editorial formatting only, in addition: an AMS \texttt{amsart} preamble that restates the article's own declarations and macros used by the retained lines, namely the environments \texttt{proposition} on separate counters, together with the standard packages those lines need.  The retained environments print in this document as Proposition 1 in order of appearance, whereas the article prints the same items with the numbering of its own full document.  The complete native source of the article as distributed in the arXiv e-print for this exact version is bundled unchanged as \texttt{sources/199105/source0.tex}.
\begin{proposition}\label{PropAct}
    \((B \times X,\cdot)\) is a \(2\)-nilpotent group if and only if the following conditions hold:
    \begin{enumerate}
        \item \(\xi\) is a derived action of \(B\) on \(X\) in the category \(\Grp\) of groups, that is:
              \begin{itemize}
                  \item[(i)] \(1 \ast x = x\);
                  \item[(ii)] \((bb') \ast x = b \ast (b' \ast x)\);
                  \item[(iii)] \(b \ast (xx') = (b \ast x)(b \ast x')\);
              \end{itemize}

        \item \((b \ast x)x^{-1} \in \bZ(X)\cap X^B\),
    \end{enumerate}
    for any \(b\), \(b' \in B\) and \(x\), \(x' \in X\), where
    \[
        X^B=\{x \in X \mid b\ast x=x, \; \forall b\in B\}
    \]
    is the set of fixed points of the action\footnote{If \(\xi\) is a derived action of \(B\) on \(X\) in \(\Grp\), then one may check that \(X^B\) is a subgroup of \(X\).}. The resulting \(2\)-nilpotent group is the \emph{semidirect product} of \(B\) and \(X\), denoted by \(B \ltimes X\).
\end{proposition}

\begin{proof}
    Suppose that \((B \times X,\cdot)\) is a \(2\)-nilpotent group. Then, \(\xi\) is a derived action in the category \(\Grp\). Moreover, for any \((b,x)\), \((b',x')\), \((c,y)\in B\times X\), the equality
    \[
        [(b,x),(b',x')] \cdot (c,y)=(c,y) \cdot [(b,x),(b',x')]
    \]
    gives
    \begin{equation}\label{eq_2nil}
        \begin{split}
            x(b \ast x')([b,b']b' \ast x^{-1}) & ([b,b'] \ast (x'^{-1}y))                                                                            \\
            =                                  & y\, c \ast \bigl(x(b \ast x')\bigl([b,b']b' \ast x^{-1}\bigr)\bigl([b,b'] \ast x'^{-1}\bigr)\bigr).
        \end{split}
    \end{equation}
    Now, for \(b'=c=1\) and \(x=1\), using the fact that \(b \ast 1=1\), we obtain
    \[
        (b \ast x')x'^{-1}y=y(b \ast x')x'^{-1},
    \]
    and therefore
    \[
        (b \ast x')x'^{-1}\in \bZ(X).
    \]
    Similarly, if we take \(b=1\) and \(x=y=1\), we get
    \[
        c \ast \bigl((b \ast x')x'^{-1}\bigr)=(b \ast x')x'^{-1},
    \]
    so that
    \[
        (b \ast x')x'^{-1}\in X^B.
    \]
    Conversely, suppose that conditions~(1) and~(2) hold. Since \(\xi\) is a derived action of~\(B\) on \(X\) in \(\Grp\), it follows that \((B \times X,\cdot)\) is a group. It remains to prove that it is \(2\)-nilpotent, that is,
    \[
        [(b,x),(b',x')] \cdot (c,y)=(c,y)\cdot [(b,x),(b',x')]
    \]
    for every \((b,x)\), \((b',x')\), \((c,y)\in B\times X\). Since both \(B\) and \(X\) are \(2\)-nilpotent, it is enough to verify that \eqref{eq_2nil} holds.

    Now, the condition \((b \ast x)x^{-1}\in X^B\) implies that
    \[
        (bb')\ast x=(b'b)\ast x
    \]
    for any \(b\), \(b'\in B\) and \(x\in X\). Indeed, one has
    \begin{align*}
        ((bb') \ast x)((b'b) \ast x)^{-1} = & ((bb') \ast x)((b'b) \ast x^{-1})                                                      \\
        =                                   & (b \ast (b' \ast x)) (b' \ast (b \ast x^{-1}))                                         \\
        =                                   & (b \ast ((b' \ast x) x^{-1})) (b \ast x) (b' \ast ((b \ast x^{-1})x))(b' \ast x^{-1} ) \\
        =                                   & (b' \ast x)x^{-1} (b \ast x) (b \ast x^{-1})x (b' \ast x^{-1} )                        \\
        =                                   & (b' \ast x)x^{-1} x (b' \ast x^{-1} )                                                  \\
        =                                   & (b' \ast x)(b' \ast x^{-1})=1.
    \end{align*}
    Therefore, \([b,b'] \ast x=x\) and \eqref{eq_2nil} reduces to
    \[
        gy = y(c\ast g),
    \]
    where
    \[
        g=x(b\ast x')(b\ast x^{-1})x'^{-1}.
    \]
    To conclude, observe that
    \[
        g=\bigl((b'\ast x)x^{-1}\bigr)^{-1}[\,b'\ast x,\; b\ast x'\,](b\ast x')x'^{-1}.
    \]
    Moreover, the commutator subgroup \([X,X]\) is fixed by the action of \(B\). Indeed, let \(b\in B\) and \(x\), \(x'\in X\). By condition~(2), there exist \(z_x\), \(z_{x'}\in \bZ(X)\) such that
    \[
        b\ast x = z_x x
        \quad\text{and}\quad
        b\ast x' = z_{x'}x'.
    \]
    It follows that
    \[
        b\ast [x,x']=[b\ast x,b\ast x']=[z_xx,z_{x'}x']=[x,x'],
    \]
    since \(z_x\) and \(z_{x'}\) are central.

    Now, since \(X^B\) is a subgroup of \(X\) and
    \[
        [X,X]\leq \bZ(X)\cap X^B,
    \]
    we deduce from the above expression for \(g\) that
    \[
        g\in \bZ(X)\cap X^B.
    \]
    Hence
    \[
        gy = yg = y(c\ast g),
    \]
    so that \eqref{eq_2nil} holds. Therefore, \((B \times X,\cdot)\) is a \(2\)-nilpotent group.
\end{proof}

\end{document}
